\subsection{Характеристика двигателя}

Будем считать, что в системе используется электродвигатель постоянного тока. Между движущим моментом двигателя $M_{\text{д}}$, управляющим сигналом $u$ и угловой скоростью ротора $\dot{\varphi}_{\text{р}}$ существует функциональная зависимость
\begin{equation}
    F(u, M_{\text{д}}, \dot{\varphi}_{\text{р}})=0.
\end{equation}

Для установившихся и медленных переходных режимов используется статическая механическая характеристика двигателя. Запишем уравнение механической характеристики:
\begin{equation}
    M_{\text{д}}
    =
    -s
    \left(
        \dot{\varphi}_{\text{р}}
        -
        \dot{\varphi}_{\text{хх}}
    \right),
\end{equation}
где $s$ --- крутизна механической характеристики двигателя.

Скорость холостого хода определяется выражением
\begin{equation}
    \dot{\varphi}_{\text{хх}}
    =
    \frac{r}{s}u,
\end{equation}
где $r$ --- коэффициент пропорциональности.

Тогда уравнение механической характеристики двигателя принимает вид
\begin{equation}
    M_{\text{д}}
    =
    -s \dot{\varphi}_{\text{р}}
    +
    r u.
\end{equation}

При исследовании переходных процессов необходимо учитывать инерционность двигателя. С учетом собственной постоянной времени двигателя $\tau$ динамическая характеристика двигателя записывается следующим образом:
\begin{equation}
    \tau \dot{M}_{\text{д}}
    +
    M_{\text{д}}
    =
    -s \dot{\varphi}_{\text{р}}
    +
    r u.
\end{equation}

Для дальнейшего анализа приведем динамическую характеристику двигателя к координате выходного звена. Умножим обе части уравнения на передаточное отношение редуктора $i$:
\begin{equation}
    i \tau \dot{M}_{\text{д}}
    +
    i M_{\text{д}}
    =
    -i s \dot{\varphi}_{\text{р}}
    +
    i r u.
\end{equation}

Так как
\begin{equation}
    M_{\text{дпр}} = i M_{\text{д}},
\end{equation}
а приведенная координата ротора определяется соотношением
\begin{equation}
    \varphi_{\text{рпр}}
    =
    \frac{\varphi_{\text{р}}}{i},
\end{equation}
то
\begin{equation}
    \dot{\varphi}_{\text{р}}
    =
    i \dot{\varphi}_{\text{рпр}}.
\end{equation}

Следовательно, динамическая характеристика двигателя в приведенных координатах имеет вид
\begin{equation}
    \tau \dot{M}_{\text{дпр}}
    +
    M_{\text{дпр}}
    =
    -s_{\text{пр}} \dot{\varphi}_{\text{рпр}}
    +
    r_{\text{пр}} u,
\end{equation}
где
\begin{equation}
    s_{\text{пр}} = i^2 s
\end{equation}
--- приведенная крутизна механической характеристики двигателя,

\begin{equation}
    r_{\text{пр}} = i r
\end{equation}
--- приведенный коэффициент пропорциональности.
